\subsection{Modelos secuenciales}
\label{sec:dl_seq_bayes}

\subsubsection*{RNN y LSTM}

Otra familia de arquitecturas fundamentales en audio son las \emph{Redes Neuronales Recurrentes} (RNNs), diseñadas para procesar secuencias temporales.  
En una RNN, la salida en cada instante $t$ depende no solo de la entrada $x_t$, sino también del estado oculto $h_{t-1}$, lo que permite modelar dependencias a lo largo del tiempo.  

El entrenamiento de RNNs estándar enfrenta el problema del \emph{gradiente desvanecido o explosivo}.  
Durante la retropropagación a través del tiempo (BPTT), los gradientes se multiplican repetidamente por derivadas de funciones de activación.  
Si estas son menores que uno, los gradientes tienden a cero (\emph{vanishing}); si son mayores que uno, crecen exponencialmente (\emph{exploding}).  
Como consecuencia, la red tiene dificultades para aprender dependencias de largo plazo de manera estable.

Para mitigar este problema se introdujeron variantes como las \emph{Long Short-Term Memory} (LSTM) y las \emph{Gated Recurrent Units} (GRU), que incorporan compuertas de control del flujo de información:
\begin{itemize}
	\item La \textbf{puerta de olvido} decide qué parte del estado anterior se descarta.  
	\item La \textbf{puerta de entrada} regula qué información nueva se incorpora a la memoria.  
	\item La \textbf{puerta de salida} determina qué fracción del estado interno influye en la salida actual.  
\end{itemize}

Gracias a este diseño, las LSTM mantienen gradientes estables en secuencias largas, permitiendo capturar dependencias temporales de distinto alcance \cite{DL_MusicReview}.  
Estas arquitecturas han sido ampliamente utilizadas en tareas de \emph{Music Information Retrieval} (MIR), como clasificación de género musical, transcripción automática o separación de fuentes, y continúan siendo un estándar en el modelado de series temporales \cite{DL_MusicReview}.


\subsection{Relevancia y diferencias frente a CNNs}

Tanto las CNNs como las RNNs trabajan con información temporal, pero de manera distinta.  
En un espectrograma, una CNN aprende patrones \textbf{locales} en bloques tiempo–frecuencia (armónicos, texturas, repeticiones), 
y puede combinarlos jerárquicamente para formar representaciones más complejas.  
Sin embargo, su ventana de contexto queda limitada por el tamaño de los filtros y la profundidad de la red.

Las RNNs, en cambio, procesan la secuencia de cuadros de manera explícita: el estado oculto transporta información entre instantes, 
lo que permite modelar dependencias \textbf{a más largo plazo} que las que cubre un kernel convolucional.  
Por eso resultan útiles cuando interesa capturar variaciones lentas o estructuras que se extienden en el tiempo 
(como la evolución del zumbido de un insecto o el fraseo en música).

En síntesis:
\begin{itemize}
	\item \textbf{CNNs}: detectan patrones locales robustos en el plano tiempo–frecuencia.  
	\item \textbf{RNNs/LSTM}: modelan cómo esos patrones se encadenan a lo largo de secuencias extensas.  
\end{itemize}

Hoy en día, ambos enfoques se consideran complementarios y a menudo se integran en arquitecturas híbridas.
